\subsection{Tensor Products of Semisimple Modules}

PROPOSITION 4. --- For \(i \in \{1,2\}\) , let \(\mathrm{A}_i\) be a K-algebra, \(\mathrm{M}_i\) a semisimple \(\mathrm{A}_i\) -module, and \(\mathrm{Z}_i\) the center of the commutant of \(\mathrm{M}_i\) . Set \(\mathrm{A} = \mathrm{A}_1 \otimes \mathrm{A}_2\) , \(\mathrm{M} = \mathrm{M}_1 \otimes \mathrm{M}_2\) , and \(\mathrm{Z} = \mathrm{Z}_1 \otimes \mathrm{Z}_2\) . We have \(\Re_{\mathrm{A}}(\mathrm{M}) = \Re(\mathrm{Z})\mathrm{M}\) . If one of the algebras \(\mathrm{Z}_1\) and \(\mathrm{Z}_2\) is separable over \(\mathrm{K}\) , in particular if the field \(\mathrm{K}\) is perfect, then the A-module \(\mathrm{M}\) is without radical.

For \(i \in \{1,2\}\) , let \(S_i\) be a simple \(A_i\) -module, \(D_i\) its commutant, and \(I(i)\) a set. We begin by treating the case when \(M_i\) is the \(A_i\) -module \(S_i^{(I(i))}\) . We identify the center \(Z_i\) of its commutant with the center of \(D_i\) . Set \(D = D_1 \otimes D_2\) .

We have \(\Re (\mathrm{D}) = \Re (\mathrm{Z})\mathrm{D}\) (Proposition 3 of VIII, p.~217) and \(\Re_{\mathrm{A}}(\mathrm{S}_1\otimes \mathrm{S}_2) =\) \(\Re (\mathrm{D})(\mathrm{S}_1\otimes \mathrm{S}_2)\) (VIII, p.~215, Corollary 2), hence \(\Re_{\mathrm{A}}(\mathrm{S}_1\otimes \mathrm{S}_2) = \Re (\mathrm{Z})(\mathrm{S}_1\otimes \mathrm{S}_2)\) . The A-module M is the direct sum of a family of A-modules isomorphic to \(\mathrm{S}_1\otimes \mathrm{S}_2\) , and the radical of the direct sum of a family of modules is the direct sum of the radicals (VIII, p.~152, Corollary 2). We therefore have \(\Re_{\mathrm{A}}(\mathrm{M}) = \Re (\mathrm{Z})\mathrm{M}\) .

Now consider the general case. For \(i \in \{1,2\}\) , denote the support of the \(A_{i}\) -module \(M_{i}\) by \(S_{M_{i}}\) . For \(\lambda \in S_{M_{i}}\) , denote the isotypical component of \(M_{i}\) of type \(\lambda\) by \(M_{i;\lambda}\) and the center of its commutant by \(Z_{i;\lambda}\) . We identify the ring \(Z_{i}\) with the product of the rings \(Z_{i;\lambda}\) for \(\lambda \in S_{M_{i}}\) . Let \(\lambda \in S_{M_{1}}\) and \(\mu \in S_{M_{2}}\) . Denote by \(i_{\lambda}: Z_{1;\lambda} \to Z_{1}\) the unique K-linear mapping such that \(pr_{\lambda} \circ i_{\lambda}\) is the identity mapping on \(Z_{1;\lambda}\) and \(pr_{\lambda'} \circ i_{\lambda}\) is the zero mapping for \(\lambda' \in S_{M_{1}} = \{\lambda\}\) . Define \(i_{\mu}: Z_{2;\mu} \to Z_{2}\) likewise. Set \(Z_{\lambda,\mu} = Z_{1;\lambda} \otimes Z_{2;\mu}\) and \(i_{\lambda,\mu} = i_{\lambda} \otimes i_{\mu}\) , and denote the mapping \(pr_{\lambda} \otimes pr_{\mu}\) from Z to \(Z_{\lambda,\mu}\) by \(\pi_{\lambda,\mu}\) . The mapping \(\pi_{\lambda,\mu}\) is a surjective ring homomorphism; we therefore have \(\pi_{\lambda,\mu}(\Re(Z)) \subset \Re(Z_{\lambda,\mu})\) (VIII, p.~155, Proposition 5, b)). Let us prove the reverse inclusion. Let z be an element of \(\Re(Z_{\lambda,\mu})\) ; since \(Z_{1;\lambda}\) and \(Z_{2;\mu}\) are fields, z is nilpotent (Theorem 3 of VIII, p.~215). We have \(i_{\lambda,\mu}(xy) = i_{\lambda,\mu}(x) \cdot i_{\lambda,\mu}(y)\) for x, y in \(Z_{\lambda,\mu}\) ; consequently, \(i_{\lambda,\mu}(z)\) is nilpotent and therefore belongs to \(\Re(Z)\) . Since \(\pi_{\lambda,\mu} \circ i_{\lambda,\mu}\) is the identity mapping on \(Z_{\lambda,\mu}\) , the element z belongs to \(\pi_{\lambda,\mu}(\Re(Z))\) . We have thus proved the equality \(\pi_{\lambda,\mu}(\Re(Z)) = \Re(Z_{\lambda,\mu})\) .

Set \(M_{\lambda,\mu}=M_{1;\lambda}\otimes M_{2;\mu}\) ; this is a submodule of M that is stable under Z. For \(z\in Z\) and \(m\in M_{\lambda,\mu}\) , we have \(zm=\pi_{\lambda,\mu}(z)m\) . Consequently, \(\Re(Z)M_{\lambda,\mu}\) is equal to \(\Re(Z_{\lambda,\mu})M_{\lambda,\mu}\) and therefore to \(\Re_{A}(M_{\lambda,\mu})\) by the isotypical case. Since the radical of a direct sum is the direct sum of the radicals (VIII, p.~152, Corollary 2) and M is the direct sum of the submodules \(M_{\lambda,\mu}\) for \((\lambda,\mu)\in\mathcal{S}_{\mathrm{M}_{1}}\times\mathcal{S}_{\mathrm{M}_{2}}\) , the equality \(\Re_{A}(M)=\Re(Z)M\) is proved. The last assertion then follows from Proposition 3 of VIII, p.~217.

Lemma 2. --- Let \(A_{1}\) and \(A_{2}\) be algebras over the commutative field K. Let \(M_{1}\) be an \(A_{1}\) -module that is finite-dimensional over K and \(M_{2}\) a \(A_{2}\) -module of finite length. Then the \(A_{1} \otimes A_{2}\) -module \(M_{1} \otimes M_{2}\) has finite length.

Set \(M = M_{1} \otimes M_{2}\) . Let \((e_{1}, \ldots, e_{n})\) be a basis of \(M_{1}\) over the field K. The mapping \((x_{1}, \ldots, x_{n}) \mapsto \sum_{i=1}^{n} e_{i} \otimes x_{i}\) is an isomorphism from the \(A_{2}\) -module \(M_{2}^{n}\) to the \(A_{2}\) -module M. Since \(M_{2}\) is an \(A_{2}\) -module of finite length, so is M. Moreover, every A-submodule of M is an \(A_{2}\) -submodule; consequently, M is an A-module of finite length.

PROPOSITION 5. --- Let \(A_{1}\) and \(A_{2}\) be algebras over the commutative field K. Let \(M_{1}\) be a semisimple \(A_{1}\) -module that is finite-dimensional over K and \(M_{2}\) a semisimple \(A_{2}\) -module. For \(i \in \{1,2\}\) , denote the commutant of the \(A_{i}\) -module \(M_{i}\) by \(D_{i}\) and the center of \(D_{i}\) by \(Z_{i}\) . Set \(A = A_{1} \otimes A_{2}\) , \(M = M_{1} \otimes M_{2}\) , \(D = D_{1} \otimes D_{2}\) , and \(Z = Z_{1} \otimes Z_{2}\) .

\begin{enumerate}
\def\labelenumi{\alph{enumi})}
\item
  The commutant of the A-module M can be identified with D, and its center is Z. If the \(A_{2}\) -module \(M_{2}\) has finite length, then the A-module M has finite length, the ring D is right and left Artinian, and the ring Z is Artinian.
\item
  The following properties are equivalent:
\end{enumerate}

\begin{enumerate}
\def\labelenumi{(\roman{enumi})}
\item
  The A-module M is semisimple.
\item
  The ring \(Z\) is isomorphic to the product of a family of commutative fields.
\item
  The ring Z is reduced.
\end{enumerate}

\begin{enumerate}
\def\labelenumi{\alph{enumi})}
\setcounter{enumi}{2}
\tightlist
\item
  The following properties are equivalent:
\end{enumerate}

\begin{enumerate}
\def\labelenumi{(\roman{enumi})}
\item
  The A-module M is isotypical and not reduced to 0.
\item
  The ring Z is a field.
\item
  The ring \(\mathbf{Z}\) is an integral domain.
\end{enumerate}

By assumption, \(M_{1}\) is finite-dimensional over K. The commutant of M can therefore be identified with D (VIII, p.~213, Lemma 1), and its center is Z (III, §4, No.~4, p.~468, Corollary). Suppose that the \(A_{2}\) -module \(M_{2}\) has finite length. The A-module M has finite length by Lemma 2. Since \(M_{2}\) is semisimple and finitely generated, the ring \(D_{2}\) is semisimple (VIII, p.~139, Proposition 6), and its center \(Z_{2}\) is the product of a finite family of commutative fields. Consequently, the \(D_{2}\) -module \((\mathrm{D}_{2})_{s}\) and the \(Z_{2}\) -module \((\mathrm{Z}_{2})_{s}\) each have finite length. Furthermore, since \(M_{1}\) is finite-dimensional over K, so are \(D_{1}\) and \(Z_{1}\) . By Lemma 2, the module \((\mathrm{D}_{1} \otimes \mathrm{D}_{2})_{s}\) has finite length; therefore, the ring \(D_{1} \otimes D_{2}\) is left Artinian. We prove likewise that the ring \(D_{1} \otimes D_{2}\) is right Artinian and that the ring \(Z_{1} \otimes Z_{2}\) is Artinian. Assertion a) follows.

Let us prove b). The center of the commutant of a semisimple module is isomorphic to the product of a family of commutative fields (VIII, p.~87, Proposition 8, a)); this proves that (i) implies (ii). The implication (ii) \(\Rightarrow\) (iii) is clear.

Suppose that the ring Z is reduced. We then have \(\Re(Z)=0\) (VIII, p.~215, Theorem 3) and \(\Re_{\mathrm{A}}(\mathrm{M})=0\) (VIII, p.~218, Proposition 4). Since the \(A_{2}\) -module \(M_{2}\) is semisimple, there exist a family \((\mathrm{S}_{i})_{i\in\mathrm{I}}\) of simple \(A_{2}\) -modules and an isomorphism from \(M_{2}\) to \(\bigoplus S_{i}\) . Consequently, the A-module M is isomorphic to \(\bigoplus M_{1}\otimes S_{i}\) . For every \(i\in I\) , the A-module \(M_{1}\otimes S_{i}\) is therefore without radical. By a), it has finite length; hence it is semisimple (VIII, p.~153, Proposition 3, b)). The A-module M is then the direct sum of a family of semisimple modules and is consequently semisimple. This proves that (iii) implies (i) and concludes the proof of b).

An A-module is isotypical and nonzero if and only if it is semisimple and the center of its commutant is a field (VIII, p.~87, Proposition 8, b)). So c) follows from b).

COROLLARY. --- If \(Z_{1}\) or \(Z_{2}\) is a separable algebra over the field K (which is, for example, the case when K is perfect), then the A-module \(M_{1} \otimes M_{2}\) is semisimple.

The rings \(Z_{1}\) and \(Z_{2}\) are isomorphic to products of fields and are therefore reduced rings. In particular, if K is perfect, they are separable algebras over K (V, §15, No.~5, p.~125, Theorem 3). By Proposition 5 of V, §15, No.~2, p.~120, the tensor product of a separable algebra and a reduced algebra is reduced. So Z is a reduced ring, and the corollary follows from Proposition 5, b).

